% Section 22.2: General properties | Part: 4 | Depends on: sec:conn, sec:lc | Status: done | Original lines: 3178-3207
%!TEX root = ../../main.tex
% REWRITTEN: sphere azimuth written \varphi instead of \phi (docs/NOTATION.md).
\subsection{General properties}
\begin{itemize}[nosep]
\item $R(X,Y,Z)=-R(Y,X,Z)$ (antisymmetric in the first two slots).
\item \textbf{First Bianchi identity.} $R(X,Y,Z)+R(Y,Z,X)+R(Z,X,Y)=\nabla_X(T(Y,Z))+\dots+T(X,[Y,Z])+\dots$ which is $0$ when $\nabla$ is torsion-free. It follows from the Jacobi identity $[X,[Y,Z]]+[Y,[Z,X]]+[Z,[X,Y]]=0$ (torsion-free case proved in Identity~\ref{id:bianchi1} of the Glossary).% REWRITTEN: glossary pointer
\item \textbf{Flat connection:} $R=0$ ($\nabla$ is \emph{flat}). Example: $\nabla^0$ on $\R^m$.
\item \textbf{Transport around loops.} If $(t,s)\mapsto\gamma(t,s)$ is a family of curves and $V$ a field along $\gamma$, then $\nabla_{\partial_s\gamma}\nabla_{\partial_t\gamma}V-\nabla_{\partial_t\gamma}\nabla_{\partial_s\gamma}V=R(\partial_s\gamma,\partial_t\gamma,V)$ (\cite[Prop.~7.5]{Lee}): the precise form of the intuition above.
% AUTHOR-NOTE: the Lee entry lists two books; check which one (and which edition) "Prop. 7.5" refers to.
\end{itemize}

\begin{example}[{the plane in polar coordinates is flat}]
Use \eqref{eq:riemanncomp} with $\Gamma^r{}_{\varphi\varphi}=-r$, $\Gamma^\varphi{}_{r\varphi}=1/r$ and compute $R^r{}_{r\varphi\varphi}$:
\[
R^r{}_{r\varphi\varphi}=\partial_r\Gamma^r{}_{\varphi\varphi}-\partial_\varphi\Gamma^r{}_{r\varphi}+\Gamma^s{}_{\varphi\varphi}\Gamma^r{}_{rs}-\Gamma^s{}_{r\varphi}\Gamma^r{}_{\varphi s}
=-1-0+0-\Gamma^\varphi{}_{r\varphi}\Gamma^r{}_{\varphi\varphi}=-1+\tfrac1r\cdot r=0 .
\]
The Christoffel symbols are nonzero but the curvature vanishes: the nonzero $\Gamma$ are only an artefact of the coordinates. \textbf{$R$ is a tensor; if a tensor vanishes in one chart it vanishes in all charts.}
\end{example}

\begin{example}[{the unit 2-sphere}]
With $\Gamma^\theta{}_{\varphi\varphi}=-\sin\theta\cos\theta$ and $\Gamma^\varphi{}_{\theta\varphi}=\cot\theta$:
\[
\begin{aligned}R^\theta{}_{\theta\varphi\varphi}&=\partial_\theta\Gamma^\theta{}_{\varphi\varphi}-\partial_\varphi\Gamma^\theta{}_{\theta\varphi}+\Gamma^s{}_{\varphi\varphi}\Gamma^\theta{}_{\theta s}-\Gamma^s{}_{\theta\varphi}\Gamma^\theta{}_{\varphi s}\\
&=-(\cos^2\theta-\sin^2\theta)-0+0-\cot\theta(-\sin\theta\cos\theta)=\sin^2\theta .\end{aligned}
\]
(The only contributing $s$ in the last term is $s=\varphi$.) So $R(\partial_\theta,\partial_\varphi)\partial_\varphi=\sin^2\theta\,\partial_\theta=g_{\varphi\varphi}\,\partial_\theta$, which is the special case of the \emph{constant-curvature formula}
\[
R(X,Y)Z=K\,\big(g(Y,Z)X-g(X,Z)Y\big),\qquad K=1\ \ (\text{unit sphere}).
\]
For the sphere of radius $a$, $K=1/a^2$. The curvature is nonzero in every chart: you cannot flatten a sphere.
\end{example}
